\section{引言}
\label{sec:1}

联邦学习允许多台边缘设备在保留本地数据的前提下协同训练全局模型。当训练节点延伸至移动终端、车载计算平台或各类传感器时，中心服务端分发全局模型参数，以及各设备回传本地计算得到的模型更新，均落在真实的物理无线链路之上。每轮通信耗费的时间，不仅取决于模型文件的大小，也直接受制于当时当地的无线传输条件。主流的联邦学习编排框架通常面向通用网络环境设计，将底层通信逻辑封装在标准套接字与传输层连接之后。对分布式算法开发而言，这种分层抽象非常省心。但若研究重点恰恰是无线环境本身，问题就会暴露出来。常规网络接口只暴露数据收发与连接状态，研究人员既无法主动设定底层无线参数，也难以在相同实验条件下，把具体的无线配置与整轮训练耗时联系起来系统比较。

在真实物理信道中，无线传输的表现取决于多项具体的收发设置。其中，调制编码方案（MCS）决定单帧的调制阶数与编码率，信道带宽决定频带宽度，符号保护间隔（GI）用来抵抗多径反射引起的时延扩展，前向纠错（FEC）在数据包层面添加校验冗余以抵御突发丢帧，发射功率则决定天线端的发射强度。常规商用 Wi-Fi 网卡为了维持连接稳定，往往由固件在后台自动调节传输速率。这种自适应机制会根据瞬时信道质量频繁升降调制阶数，导致每次发送的物理参数自行漂移。对受控实验而言，这就带来了麻烦。实验人员既无法锁定一组指定的无线参数开展严谨测试，也难以在不同轮次间重现完全相同的传输设置。一旦某个客户端回传滞后，或者整轮耗时大幅拉长，研究人员仅凭应用层暴露的耗时数字，很难弄清通信表现究竟与哪些无线参数直接相关。

针对这一局限，本文提出基于消费级 802.11 硬件，把联邦学习中的真实无线链路转变为可主动配置条件且可记录接收信息的实验对象。平台的设计目标，是支持研究人员在固定模型结构、数据集与计算条件的前提下，主动设定不同的无线参数组合，进而在完整训练流程中定量对比模型分发吞吐、客户端更新提交进展以及整轮耗时。在技术途径上，平台借助消费级网卡的监听模式（monitor 模式）与原始帧注入绕过商用固件的控制逻辑。网卡开启监听模式后脱离常规接入点的关联流程，直接捕获空中实际到达的无线帧，并记录随帧附带的物理层元数据，例如实际接收到的调制编码方案索引与信号强度。在模型下行分发时由客户端记录接收信息，在上行更新回传时由服务端记录。原始帧注入则赋予发送端绕过系统默认协商机制的能力，直接按照实验预设的参数构造并送出 802.11 物理帧。各项无线条件的配置分别落在不同的系统层次。工作信道、带宽与发射功率通过驱动与内核接口直接下发给射频芯片。调制编码方案与保护间隔由发送端在构造单帧时写入发射描述信息，由物理层逐帧采纳。前向纠错则运行在链路控制层，通过通用处理器对数据分组计算块级校验冗余，与底层射频配置分属不同层级。

消费级 802.11 硬件为这类受控实验带来了切实的设备可获得性。市售的双频 USB 无线网卡单价在百元级，依托免许可频段工作，不需要申请专门的无线电执照，研究团队很容易备齐多节点实验所需的外设。不过，设备易得并不等同于系统开箱即用。这套方案依赖特定网卡芯片与开源驱动补丁来支持监听和帧注入，实验人员依然需要在多台边缘节点上逐一配置系统环境、编译专用驱动并联调收发控制逻辑。此外，实验参数可配置，并不代表真实世界的复杂电磁环境能够被完全掌控或精确复制。通信对端记录的帧级元数据，也无法替代专用射频仪表的连续信道估计。平台的作用，在于为研究人员提供一套低成本且可行的实验工具，在真实硬件上把物理无线配置与端到端联邦学习流程结合起来，支持有依据的对比测试与通信特征分析。
